\documentclass{article}
\usepackage[utf8]{inputenc}
\usepackage[T1]{fontenc}
\usepackage[a4paper,total={6in,8in}]{geometry}
\usepackage{amsmath,amssymb,amsfonts}
\usepackage{graphicx}
\usepackage{hyperref}
\usepackage{enumitem}
\usepackage{listings}
\usepackage{xcolor}
\DeclareMathOperator{\Var}{Var}
\DeclareMathOperator{\Cov}{Cov}
\providecommand{\tightlist}{\setlength{\itemsep}{0pt}\setlength{\parskip}{0pt}}
\lstset{basicstyle=\ttfamily\scriptsize,breaklines=true,columns=fullflexible,
  keepspaces=true,showstringspaces=false,frame=single}
\hypersetup{colorlinks=true,urlcolor=blue,linkcolor=blue}
\setlength{\emergencystretch}{2em}

\title{Math 126 Problem 8: Sphere in a Box}
\author{} % Put your name here.
\date{Spring 2027}
\begin{document}
\maketitle

\hypertarget{submission-instructions}{%
\subsection*{Submission Instructions}\label{submission-instructions}}

Please upload your work to Gradescope by \textbf{11:59pm on Monday, February 22, 2027 (Pacific time).}

\begin{itemize}
\item
  You are encouraged to discuss ideas and work with your classmates. However, you \emph{must acknowledge} your collaborators at the top of each solution on which you collaborated with others and you \emph{must write} your solutions independently.
\item
  Your solutions to theory questions must be written legibly, or typeset in LaTeX or markdown.
\item
  I recommend that you write your solutions in a Python notebook using Google Colab.
\item
  You should submit your solutions as a \emph{single PDF} via the assignment on Gradescope.
\end{itemize}

\textbf{Grading:} The point of the problem set is for \emph{you} to learn. To this end, I hope to disincentivize the use of LLMs by \emph{not} grading your work for correctness. Instead, you will grade your own work by comparing it to my solutions. This self-grade is due the Friday \emph{after} the problem set is due, also on Gradescope.

\noindent Original notebook: \href{https://colab.research.google.com/drive/1HY7qkCvXM856OxkEsTJdbbgv-AZJvyqh?usp=sharing}{Google Colab}.

% Source notebook cell 2
\hypertarget{sphere-in-a-box}{%
\subsection*{Sphere in a Box}\label{sphere-in-a-box}}

Consider a \(d\)-dimensional hypercube centered at the origin with side length \(4\) (so it spans \([-2, 2]\) in each dimension). Place \(2^d\) unit spheres (radius \(r=1\)) into the corners of the hypercube. Each sphere is centered at a coordinate \((\pm 1, \pm 1, \dots, \pm 1)\).

Now, place a smaller central sphere at the origin so that it exactly touches all \(2^d\) corner spheres. Using the Pythagorean theorem, calculate the distance from the origin to the center of a corner sphere. Deduce the radius \(R_d\) of this central sphere.

Here is the weird part. We know the distance from the origin to the face of the hypercube is 2 (since the side length is 4). Find the dimension \(d\) where \(R_d = 2\). (At this point, the central sphere is tangent to the outside walls of the box). What happens when \(d > 9\)?

Describe the shape of the central sphere relative to the hypercube for extremely large \(d\). How can a sphere enclosed by other spheres poke outside the box containing them?

\end{document}
